% Complete kernel construction, source conditional laws and coordinate examples.

\begin{annexitem}{annex:p15.kernel-expectation}{thm:p15-kernel-expectation}{Measurable integral approximation with one measure at each argument}
\begin{proof}
For an indicator \(f=\mathbf1_A\), the kernel integral is \(K(y,A)\),
measurable by the kernel definition.
Finite nonnegative linear combinations give measurability for simple
functions. For arbitrary nonnegative \(f\), choose fixed simple
approximations \(f_n\uparrow f\) on \(S\).
At every parameter \(y\), monotone convergence for the probability
measure \(K(y,\cdot)\) gives
\[
                    \int f\,dK(y,\cdot)
                        =\lim_n\int f_n\,dK(y,\cdot).
\]
The right side is a measurable increasing limit.

If \(K\) is a regular conditional law of \(X\), its defining
conditional identity yields, for every \(B\in\mathcal G\),
\[
                 \int_B K(\omega,A)\,dP(\omega)
                         =P(B\cap\{X\in A\}).
\]
Finite sums give the corresponding identity for each \(f_n\).
Monotone convergence, first inside the kernel and then on the
parameter probability space, gives
\[
          \int_B\left(\int f\,dK(\omega,\cdot)\right)dP(\omega)
                              =E(f(X)\mathbf1_B).
\]
The nonnegative conditional-expectation uniqueness theorem proves
the asserted version formula, including infinite values.

For signed \(f\) with \(E|f(X)|<\infty\), apply this to \(|f|\).
The absolute kernel integral is finite outside a
\(\mathcal G\)-measurable null set.
There the two part integrals are finite; their difference is
measurable, and its absolute value is bounded by the absolute
kernel integral. Assign zero on the exceptional set.
The resulting variable is integrable. Subtracting the two
nonnegative event identities proves it is the required version.
The full-measure domain here may depend on \(f\), even though
the probability kernel itself is defined on every argument.
\end{proof}
\end{annexitem}

\begin{annexitem}{annex:p15.real-conditional-law}{thm:p15-real-conditional-law}{All CDF conditions and the generating-class construction}
\begin{proof}
\paragraph{A single good set for the rational family.}
For every \(q\in\mathbb Q\), choose a \(\mathcal G\)-measurable
version \(g_q\) of \(P(W\leq q\mid\mathcal G)\).
Clip it to \([0,1]\), which changes it only on a conditioning-null
set by the already proved conditional bounds.
For rational \(q<r\), conditional order gives \(g_q\leq g_r\)
almost everywhere. Conditional dominated convergence gives
\[
 g_{q+1/n}\longrightarrow g_q,\qquad
 g_n\longrightarrow1,\qquad g_{-n}\longrightarrow0
                         \quad\text{almost everywhere}.
\]
The first assertion follows because the corresponding indicators
decrease to \(\mathbf1_{\{W\leq q\}}\); the other two follow
from the finiteness of the real variable \(W\).
There are countably many rational pairs and rational \(q\)'s.
Let \(\Omega_0\) be the set on which all these monotonicities
and limits hold. It is \(\mathcal G\)-measurable and has probability one.

\paragraph{The CDF for every good argument.}
For \(\omega\in\Omega_0\), define
\(F_\omega(t)=\inf_{q\in\mathbb Q,\ q>t}g_q(\omega)\).
The set of indices is nonempty, and all values lie in \([0,1]\).
As \(t\) increases, the index set decreases, so \(F_\omega\)
is nondecreasing. If \(t_n\downarrow t\), its decreasing limit
is at least \(F_\omega(t)\).
For any rational \(q>t\), eventually \(t_n<q\), so
\(F_\omega(t_n)\leq g_q(\omega)\).
Taking the infimum over such \(q\) gives the reverse inequality,
proving right continuity.

At a rational \(q\), monotonicity gives \(F_\omega(q)\geq g_q(\omega)\),
while the sequence \(q+1/n\) gives
\(F_\omega(q)\leq\lim_n g_{q+1/n}(\omega)=g_q(\omega)\).
Thus they agree on the rationals. In particular
\(F_\omega(n)\to1\) and \(F_\omega(-n)\to0\).
Monotonicity then gives the two required limits at infinite endpoints.
So \(F_\omega\) is a genuine cumulative distribution function.
For \(\omega\notin\Omega_0\), use \(F_\omega(t)=\mathbf1_{\{t\geq0\}}\).
For every fixed real \(t\), this piecewise formula is
\(\mathcal G\)-measurable in \(\omega\).

For any real \(t\), choose rationals \(q_n>t\) decreasing to \(t\).
Monotonicity shows \(g_{q_n}(\omega)\downarrow F_\omega(t)\)
on \(\Omega_0\): every rational \(q>t\) eventually exceeds \(q_n\).
Conditional dominated convergence identifies that limit almost
everywhere with \(P(W\leq t\mid\mathcal G)\).
This assertion is for each fixed \(t\), which is all that its
conditional version identity requires.

\paragraph{A measurable family of probability measures.}
For each \(\omega\), the Stieltjes construction gives the unique
probability measure \(\nu_\omega\) on \(\mathbb R\) with CDF \(F_\omega\).
We use the following general fact about this construction.
On any measurable parameter space, a family of CDFs whose value at
each fixed \(t\) is measurable gives a probability kernel.
Indeed, the class of Borel \(A\) for which the corresponding
measure of \(A\) is measurable contains all half-lines.
It contains \(\mathbb R\), is closed under complements by subtraction
from one, and under countable disjoint unions by the pointwise
sum of measure values. It is a lambda-system containing the
half-line pi-system, so contains every Borel set.
This proves the general fact and, in particular, that
\(\omega\mapsto\nu_\omega(A)\) is \(\mathcal G\)-measurable.

Now restrict the class further by requiring the event identity
\[
                   \int_B\nu_\omega(A)\,dP(\omega)
                                 =P(B\cap\{W\in A\})
                                       \qquad(B\in\mathcal G).
\]
It too is a lambda-system: the whole-space identity is \(P(B)=P(B)\);
complements subtract from it; and disjoint unions use monotone
convergence and probability additivity.
The preceding CDF version calculation puts all half-lines in this
class. The pi-lambda theorem proves the identity for every Borel \(A\).
Together these facts give the required regular conditional law.

\paragraph{One null set for uniqueness as measures.}
If \(K,L\) are two probability kernels satisfying these conditional
identities, their values on each rational half-line are versions of
the same scalar conditional probability. They agree outside a
\(\mathcal G\)-null set for that half-line.
Remove the countable union of those sets. At every remaining
argument the two probability measures agree on all rational half-lines
and on the whole space. This is a generating pi-system, so probability
uniqueness gives agreement on every Borel event at once.
\end{proof}
\end{annexitem}

\begin{annexitem}{annex:p15.standard-conditional-law}{thm:p15-standard-conditional-law}{Borel-image concentration and transport back}
\begin{proof}
Let \(W=\phi(X)\), a real random variable, and take its regular
conditional law \(\nu_\omega\) from the preceding theorem.
Since \(E=\phi(S)\) is Borel and \(W\in E\) everywhere,
\(\nu_\omega(E)\) is a version of the constant conditional
probability one. Thus
\(\Omega_1=\{\omega:\nu_\omega(E)=1\}\) belongs to \(\mathcal G\)
and has probability one.
Because \(X\) exists on a probability space, \(S\) is nonempty;
choose a fixed point \(s_0\in S\).

For \(A\in\mathcal S\), the image \(\phi(A)\) is Borel in \(E\),
since it is the inverse image of \(A\) under the measurable inverse
\(\phi^{-1}:E\to S\).
The Borel nature of \(E\) then makes \(\phi(A)\) Borel in \(\mathbb R\).
Define
\[
 K(\omega,A)=
 \begin{cases}
  \nu_\omega(\phi(A)),&\omega\in\Omega_1,\\
  \mathbf1_A(s_0),&\omega\notin\Omega_1.
 \end{cases}
\]
For each \(A\), this is \(\mathcal G\)-measurable.
At a good argument it is a probability measure on \(S\):
the bijection \(\phi\) preserves disjoint unions and the whole
image has \(\nu_\omega\)-mass one.
At any other argument it is the point mass at \(s_0\).
The conditional identity for \(\phi(A)\) and
\(\{W\in\phi(A)\}=\{X\in A\}\) proves that \(K\) is the
required regular conditional law.

For uniqueness, the sets
\(\phi^{-1}(E\cap(-\infty,q])\), \(q\in\mathbb Q\), together
with \(S\), form a countable generating pi-system for \(\mathcal S\).
Two kernels have almost-everywhere equal values on each of these
events. Outside their countable exceptional union, probability
uniqueness gives equality of the measures on all of \(\mathcal S\).
The previously proved coding theorem supplies this hypothesis for
the stated Borel subsets of complete separable metric spaces.
\end{proof}
\end{annexitem}

\begin{annexitem}{annex:p15.value-conditional-law}{thm:p15-value-conditional-law}{A kernel on every observation argument}
\begin{proof}
Choose a Borel real coding \(\phi:S\to E\), and put \(W=\phi(X)\).
For each rational \(q\), scalar conditional expectation gives
a \(\sigma(Y)\)-measurable version of \(P(W\leq q\mid\sigma(Y))\).
The proved factorization theorem represents it as \(H_q(Y)\)
for a \(\mathcal T\)-measurable real \(H_q\).
Clip every \(H_q\) to \([0,1]\), preserving the conditional versions.

Let \(T_0\) be the set of \(y\)'s on which all the countable
conditions
\[
 H_q(y)\leq H_r(y)\ (q<r),\quad
 H_{q+1/n}(y)\to H_q(y),\quad
 H_n(y)\to1,\quad H_{-n}(y)\to0
\]
hold, with the first two indexed by rationals.
It is \(\mathcal T\)-measurable. The conditional order and
convergence arguments in the real-law proof apply to \(H_q(Y)\),
so \(P(Y\in T_0)=\mu_Y(T_0)=1\).
For \(y\in T_0\), put
\(F_y(t)=\inf_{q>t,\ q\in\mathbb Q}H_q(y)\);
outside \(T_0\), use the CDF of \(\delta_0\).
The fully proved CDF argument gives a CDF at every \(y\).
Its values at fixed \(t\) are \(\mathcal T\)-measurable.
The general measurable-CDF construction proved above therefore
gives a probability kernel \(y\mapsto\nu_y\) on the real line.

For rational \(q\), the value \(\nu_Y((-\infty,q])=H_q(Y)\)
almost everywhere is the required scalar conditional version.
The class of Borel \(B\) for which \(\nu_Y(B)\) has that
conditional version property is a lambda-system, by the
whole-space, complement and disjoint-sum calculations from the
real-law proof. Rational half-lines generate, so the property
holds for every real Borel set.
In particular \(\nu_Y(E)=1\) almost everywhere.
The measurable set \(T_1=\{y:\nu_y(E)=1\}\) has \(\mu_Y\)-mass one.

Choose \(s_0\in S\), and define \(K(y,A)=\nu_y(\phi(A))\) on
\(T_1\) and \(K(y,A)=\mathbf1_A(s_0)\) elsewhere.
The same Borel-image and transport arguments as above make this
a probability kernel at every \(y\), measurable on \((T,\mathcal T)\).
Its composition with \(Y\) is a regular conditional law of \(X\)
given \(\sigma(Y)\). Applying the defining identity to
\(\{Y\in B\}\) gives
\[
                    P(X\in A,Y\in B)=\int_BK(y,A)\,d\mu_Y(y).
\]

For uniqueness, compare two kernels on the countable target
pi-system obtained from \(\phi\)'s rational half-lines.
Their compositions with \(Y\) are equal almost everywhere on
each such event, which means the kernels agree on it outside
a \(\mu_Y\)-null measurable subset of \(T\).
Remove their countable union and apply probability uniqueness
at each remaining \(y\). This gives simultaneous agreement on
all target events. No coding or measurable-image assumption on
the conditioning space \(T\) was used.
\end{proof}
\end{annexitem}

\begin{annexitem}{annex:p15.density-conditional-kernel}{thm:p15-density-conditional-kernel}{Normalization on the exact marginal support}
\begin{proof}
Tonelli makes \(m\) measurable and gives \(\int m(y)\,dy=1\).
Thus \(\{m=\infty\}\) is Lebesgue-null. The \(Y\)-law has density
\(m\), so both that set and \(\{m=0\}\) have zero \(Y\)-probability.
The Borel set \(D\) therefore has full \(Y\)-probability.
For each Borel \(A\), Tonelli applied to \(\mathbf1_A(x)f(x,y)\)
makes its section integral measurable.
The ratio on \(D\) and point-mass assignment on \(D^c\) consequently
give a Borel function of \(y\).

At each \(y\in D\), the numerator is a measure in \(A\), constructed
by a nonnegative density, with finite positive total mass \(m(y)\).
Its normalization is a probability measure.
At every other \(y\) the given point mass is also a probability measure.
For Borel \(A,B\), Tonelli and marginal integration yield
\[
 \begin{aligned}
 E(K(Y,A)\mathbf1_{\{Y\in B\}})
   &=\int_{B\cap D}\int_A f(x,y)\,dx\,dy\\
   &=P(X\in A,Y\in B).
 \end{aligned}
\]
The omitted observation set has zero marginal and joint probability.
Every event in \(\sigma(Y)\) is a Borel preimage under \(Y\);
hence this is the required conditional event identity.
\end{proof}
\end{annexitem}

\begin{annexitem}{annex:p15.kernel-joint-measure}{thm:p15-kernel-joint-measure}{Measurable sections and a complete kernel integration formula}
\begin{proof}
For product measurable \(E\subseteq T\times S\), let
\(E_y=\{x:(y,x)\in E\}\).
These sections are measurable: the sets with that property form
a sigma-algebra containing all rectangles.
Let \(\mathcal D\) consist of the product measurable sets for which
\(y\mapsto K(y,E_y)\) is measurable.
It contains rectangles because
\(K(y,(B\times A)_y)=\mathbf1_B(y)K(y,A)\).
It contains the whole space and is closed under complements by
subtraction from one. For disjoint \(E_n\), their sections are
disjoint and the value on their union is
\(\sum_nK(y,(E_n)_y)\), a measurable increasing sum.
Thus \(\mathcal D\) is a lambda-system containing the generating
rectangle pi-system, and includes every product measurable set.

Define \(Q(E)=\int_TK(y,E_y)\,d\mu(y)\).
It is nonnegative, has empty value zero and total mass one.
For disjoint unions the pointwise section identity and monotone
convergence give \(Q(\bigcup_nE_n)=\sum_nQ(E_n)\).
Thus \(Q\) is a probability measure with the stated rectangle values.
Any other such probability agrees on the rectangle pi-system and
is equal to \(Q\) by the probability uniqueness theorem.

For \(h=\mathbf1_E\), the integral formula is the definition of \(Q\).
Finite nonnegative sums prove it for simple \(h\).
For general nonnegative \(h\), take fixed simple approximations
increasing to it. Monotone convergence inside each probability
\(K(y,\cdot)\), and then for \(\mu\) and \(Q\), gives both
parameter measurability and the full integral formula.
For signed integrable \(h\), the formula for \(|h|\) makes its
absolute inner integral finite outside a measurable \(\mu\)-null set.
There subtract the two finite part integrals and set the result
to zero on the exceptional set.
The resulting measurable function has absolute value bounded by
the absolute inner integral, so is integrable.
Subtracting the two nonnegative formulas gives the signed identity.

For the conditional kernel with \(\mu=\mu_Y\), the rectangle values
coincide with the actual law of \((Y,X)\), hence uniqueness identifies
that law with \(Q\).
Apply the signed formula also to \(\mathbf1_B(y)h(y,x)\) for every
\(B\in\mathcal T\). It gives exactly the defining conditional
event identity for the inner-integral function evaluated at \(Y\).
This function is measurable from \(Y\) and integrable, proving the
last assertion.
\end{proof}
\end{annexitem}

\begin{annexitem}{annex:p15.conditional-holder}{thm:p15-conditional-holder}{The ordinary inequality on the conditional pair distribution}
\begin{proof}
The pair \((U,V)\) is a measurable random element of Euclidean
\(\mathbb R^2\), a complete separable metric space.
Its Borel real coding and the proved conditional-law theorem give
a probability kernel \(K_\omega\) for the pair given \(\mathcal G\).
Ordinary H\"older first proves \(E|UV|<\infty\).
Kernel integration identifies the three conditional quantities
with versions of
\[
        \int|uv|\,dK_\omega,\qquad
        \int|u|^p\,dK_\omega,\qquad
        \int|v|^q\,dK_\omega .
\]
All three integrals are finite outside one conditioning-null set,
since their expectations are finite.
At each remaining \(\omega\), apply ordinary H\"older on the
probability space \((\mathbb R^2,K_\omega)\).
Its inequality, including the possible zero-moment cases, is
exactly the claimed conditional inequality for these versions.
\end{proof}
\end{annexitem}

\begin{annexitem}{annex:p15.conditioning-coordinate-examples}{thm:p15-conditioning-coordinate-examples}{Spherical coordinates and two shrinking-event families}
\begin{proof}
\paragraph{The spherical model.}
For the parametrization \(r\), direct differentiation gives
\[
 \begin{aligned}
 r_\theta&=(-\cos\phi\sin\theta,\ \cos\phi\cos\theta,\ 0),\\
 r_\phi&=(-\sin\phi\cos\theta,\ -\sin\phi\sin\theta,\ \cos\phi).
 \end{aligned}
\]
The length of their cross product is \(\cos\phi\) on the latitude
interval. Thus the usual parametrized area element is
\(\cos\phi\,d\theta\,d\phi\); integrating it gives
\((2\pi)\int_{-\pi/2}^{\pi/2}\cos\phi\,d\phi=4\pi\).
This explains the normalized coordinate density specifying the model.
It factors as \((1/(2\pi))(\cos\phi/2)\), with each factor
normalized on its declared interval.
The joint-density independence theorem gives the original
marginals and independence.

On \(\Theta<\pi\), the new coordinates are \((\Theta,\Phi)\),
with density \(\cos\psi/(4\pi)\) on the central angle interval.
On \(\Theta\geq\pi,\Phi\geq0\), they are
\((\Theta-\pi,\pi-\Phi)\), a translation and reflection with
unit absolute scale in both coordinates.
Here \(\cos\Phi=-\cos\Psi=|\cos\Psi|\).
On \(\Theta\geq\pi,\Phi<0\), they are
\((\Theta-\pi,-\pi-\Phi)\), again with unit absolute scale
and the same equality of absolute cosines.
One-dimensional affine substitution in these regions gives
the joint density
\[
            \frac{|\cos\psi|}{4\pi}
                \quad(0\leq\theta'<\pi,\ -\pi<\psi\leq\pi).
\]
Boundary overlaps and coordinate duplications have zero probability.
The elementary sine and cosine identities also give
\(r(\Theta',\Psi)=r(\Theta,\Phi)\) in all three cases.
The described point on the sphere is the same.
Since \(\int_{-\pi}^{\pi}|\cos\psi|\,d\psi=4\), the new density
factors as \((1/\pi)(|\cos\psi|/4)\).
This proves both new marginals and their independence.
The density kernel, or independence directly, gives the stated
nonuniform conditional angle law.

\paragraph{Changing the uniform-square observation.}
For \(Z=Y/(1+X)\), fix \(x\in[0,1]\) and substitute
\(y=(1+x)z\) in the uniform integral.
The proved affine Lebesgue rule and Tonelli give the joint
density representative
\[
 f_{X,Z}(x,z)=(1+x)\mathbf1_{[0,1]}(x)
                         \mathbf1_{[0,\,1/(1+x)]}(z).
\]
The equality is valid on every product measurable event by
applying the same substitution to its nonnegative indicator.
For \(0\leq z<1/2\), every \(x\in[0,1]\) is in its support,
and the marginal density is \(\int_0^1(1+x)\,dx=3/2\).
The density-kernel formula therefore gives
\((2/3)(1+x)\mathbf1_{[0,1]}(x)\), with mean
\[
                       \frac23\int_0^1(x+x^2)\,dx
                           =\frac23\left(\frac12+\frac13\right)=\frac59.
\]
The original unit-square density gives uniform conditional
\(X\)-density and mean \(1/2\) at every \(y\in[0,1]\), including
the chosen density version at zero.

For \(0<\varepsilon<1/2\), independence gives
\[
 P(X\in A,Y<\varepsilon)=\varepsilon\lambda(A\cap[0,1]),\qquad
 P(Y<\varepsilon)=\varepsilon.
\]
The other event is \(Y<\varepsilon(1+X)\), whose upper bound is
less than one throughout the unit square. Thus
\[
 P(X\in A,Z<\varepsilon)
       =\varepsilon\int_{A\cap[0,1]}(1+x)\,dx,\qquad
 P(Z<\varepsilon)=3\varepsilon/2 .
\]
These are positive-probability conditioning events. Dividing
their joint probabilities by their probabilities proves exactly
the two conditional densities stated in the theorem.
Finally \(1+X>0\), so \(\{Y=0\}=\{Z=0\}\), and this is the
intersection of either shrinking family.
It is null. The different calculations use different observed
coordinates and neighborhoods, and make no claim that a
ratio conditional on this null event exists.
\end{proof}
\end{annexitem}
